
\documentclass[11pt]{article}

\usepackage[margin=1in]{geometry}
\usepackage{amsmath,amssymb}
\usepackage{listings}
\usepackage{xcolor}
\usepackage{hyperref}

\definecolor{backcolour}{rgb}{0.95,0.95,0.92}

\lstset{
    backgroundcolor=\color{backcolour},
    basicstyle=\ttfamily\small,
    breaklines=true,
    frame=single
}

\title{CPSC 250 \\ Dictionaries}
\author{Edward Brash}
\date{}

\begin{document}

\maketitle

\section{Introduction}

Dictionaries are one of the most powerful and useful data structures in Python.

Many computer scientists use dictionaries constantly because they provide an efficient way to map one piece of information to another.

Students often struggle with dictionaries initially, but they are extremely useful once understood properly.

\section{When Should You Use a Dictionary?}

A dictionary is useful whenever there exists a natural one-to-one mapping between two collections of information.

For example:

\begin{itemize}
    \item names $\rightarrow$ ages
    \item student IDs $\rightarrow$ GPAs
    \item usernames $\rightarrow$ passwords
    \item countries $\rightarrow$ capitals
\end{itemize}

If there is a relationship between two kinds of data, a dictionary may be an appropriate structure.

\section{Example Using Two Lists}

Suppose we initially store information in two separate lists:

\begin{lstlisting}[language=Python]
names = ['Bob', 'Alice', 'Jane', 'Fred']
ages = [18, 31, 16, 47]
\end{lstlisting}

The problem is that the relationship between the lists must be maintained manually.

For example:

\begin{itemize}
    \item names[0] corresponds to ages[0]
    \item names[1] corresponds to ages[1]
\end{itemize}

This becomes cumbersome and error-prone.

\section{Replacing with a Dictionary}

Instead, we can directly store the relationship:

\begin{lstlisting}[language=Python]
people = {
    'Bob': 18,
    'Alice': 31,
    'Jane': 16,
    'Fred': 47
}
\end{lstlisting}

This preserves the correlation between the data naturally.

\section{Accessing Dictionary Values}

Dictionary values are accessed using keys:

\begin{lstlisting}[language=Python]
print(people['Bob'])
\end{lstlisting}

Output:

\begin{lstlisting}
18
\end{lstlisting}

\section{Important Terminology}

In a dictionary:

\begin{lstlisting}
'Bob' : 18
\end{lstlisting}

\begin{itemize}
    \item \texttt{'Bob'} is called the \textbf{key}
    \item \texttt{18} is called the \textbf{value}
\end{itemize}

Dictionaries store \textbf{key/value pairs}.

\section{Keys and Values Can Have Different Types}

Keys and values do not need to have the same type.

Example:

\begin{lstlisting}[language=Python]
friends = {
    'Bob': 'Alice',
    'Jane': 'Heather',
    'Fred': 'John'
}
\end{lstlisting}

Another example:

\begin{lstlisting}[language=Python]
finish_order = {
    1: 'Bob',
    2: 'Fred',
    3: 'Bill'
}
\end{lstlisting}

\begin{lstlisting}[language=Python]
print(finish_order[2])
\end{lstlisting}

Output:

\begin{lstlisting}
Fred
\end{lstlisting}

Important:

The value inside square brackets is a \textbf{key}, not an index.

\section{Useful Dictionary Methods}

\begin{itemize}
    \item \verb|my_dict.keys()| --- return all keys

    \item \verb|my_dict.values()| --- return all values

    \item \verb|my_dict.items()| --- return all key/value pairs

    \item \verb|my_dict.get(key)| --- safely retrieve a value

    \item \verb|my_dict.update(other_dict)| --- merge dictionaries

    \item \verb|my_dict.pop(key)| --- remove a key/value pair
\end{itemize}

\section{Mutability}

Dictionaries are mutable structures.

This means:

\begin{itemize}
    \item values can be changed,
    \item new entries can be added,
    \item entries can be removed.
\end{itemize}

Example:

\begin{lstlisting}[language=Python]
people['Bob'] = 19
people['Mary'] = 27
\end{lstlisting}

\section{Summary of Python Data Structures}

\begin{center}
\begin{tabular}{|c|c|c|}
\hline
Type & Description & Mutable? \\
\hline
int & Numeric integer & Yes* \\
float & Floating point number & Yes* \\
string & Sequence of characters & No \\
list & Ordered sequence & Yes \\
tuple & Ordered immutable sequence & No \\
set & Unordered unique collection & Yes \\
dict & Key/value mapping & Yes \\
\hline
\end{tabular}
\end{center}

\vspace{0.5cm}

\noindent
*Python integers and floats are technically objects. Operations create new objects rather than modifying existing ones in place.

\section{Key Takeaways}

\begin{itemize}
    \item Dictionaries store mappings between keys and values.
    \item They are ideal for representing related information.
    \item Keys are used for lookup instead of numeric indexing.
    \item Dictionaries are mutable and highly flexible.
    \item Dictionaries are among the most useful data structures in Python.
\end{itemize}

\end{document}
